% from ega4-9-fr.tex
\oldpage[IV-3]{57}
\begin{remarks}[9.2.2]
\label{IV.9.2.2-fr}
\begin{enumerate}[label=(\roman*)]
  \item Il s'agit là bien entendu d'une convention de langage de nature
  métamathématique et non d'une définition mathématique proprement dite. On a
  des «~définitions~» analogues pour des relations entre $k$, un ou plusieurs
  $k$-préschémas algébriques, des $k$-morphismes entre ces préschémas, des
  Modules cohérents sur ces préschémas ou des homomorphismes entre ces Modules.

  \item Nous aurons aussi à considérer des relations où figurent des
  \emph{parties constructibles} de préschémas. Par exemple, soit
  $\mathbf P(X,Z,k)$ une relation~; nous dirons (par abus de langage) que
  $\mathbf P$ est une propriété constructible (resp. ind-constructible) de la
  partie constructible $Z$ de $X$ si les deux conditions suivantes sont
  vérifiées~:
  \begin{enumerate}[label=$\arabic*^{\mathrm o}$]
    \item Si $k$ est un corps, $X$ un préschéma algébrique sur $k$, $Z$ une
    \emph{partie constructible de $X$}, $k'$ une extension de $k$, alors, pour
    que $\mathbf P(X,Z,k)$ soit vraie, il faut et il suffit que
    $\mathbf P(X_{(k')},p^{-1}(Z),k')$ soit vraie
    ($p:X_{(k')}\to X$ étant la projection canonique).
    \item Soient $S$ un préschéma intègre noethérien, de point générique
    $\eta$, $u:X\to S$ un morphisme de type fini, $Z$ une
    \emph{partie constructible de $X$}. Pour tout $s\in S$, posons
    $X_s=u^{-1}(s)$, $Z_s=Z\cap X_s$. Soit $E$ l'ensemble des $s\in S$ tels
    que $\mathbf P(X_s,Z_s,k(s))$ soit vraie. Alors l'un des ensembles $E$,
    $S-E$ (resp. l'ensemble $E$) contient un ouvert non vide (resp. contient
    un ouvert non vide s'il contient $\eta$).
  \end{enumerate}

  On notera que dans la condition $2^{\mathrm o}$ il faut supposer que $Z$ est
  une partie constructible de $X$, et non seulement que $Z_s$ est une partie
  constructible de $X_s$ pour tout $s$~; la première de ces deux propriétés
  entraîne la seconde (1.8.2), mais non réciproquement.

  \item Si $\mathbf P$ est une propriété constructible, il est clair qu'il en
  est de même de «~non $\mathbf P$~». Si $\mathbf P$, $\mathbf Q$ sont deux
  propriétés constructibles (resp. ind-constructibles), il en est de même des
  propriétés «~$\mathbf P$ ou $\mathbf Q$~» et «~$\mathbf P$ et
  $\mathbf Q$~»~; en effet, si, sous les hypothèses de (9.2.1, (ii)), $E$,
  $E'$ sont deux parties de $S$ et si $E$ contient un ouvert non vide, il en
  est de même de $E\cup E'$, et si $S-E$ et $S-E'$ contiennent chacun un
  ouvert non vide, il en est de même de
  $S-(E\cup E')=(S-E)\cap(S-E')$.

  \item Soit $\mathbf P(X,\mathcal F,k)$ une relation vérifiant la condition
  (9.2.1, (i))~; soient $S$ un préschéma, $u:X\to S$ un morphisme de type
  fini~; avec les notations de (9.2.1, (ii)), soit $E$ l'ensemble des $s\in S$
  tels que $\mathbf P(X_s,\mathcal F_s,k(s))$ soit vraie. Soit d'autre part
  $g:S'\to S$ un morphisme quelconque, et posons
  $X'=X_{(S')}$, $\mathcal F'=\mathcal F\otimes_{\mathcal O_S}\mathcal O_{S'}$~;
  alors il résulte de la transitivité des fibres (I, 3.6.4) et de la condition
  (9.2.1, (i)) que \emph{l'ensemble des $s'\in S'$ tels que
  $\mathbf P(X'_{s'},\mathcal F'_{s'},k(s'))$ soit vraie est égal à
  $g^{-1}(E)$}. Cela s'étend aussitôt au cas où il figure dans $\mathbf P$
  plusieurs préschémas, Modules, morphismes de préschémas ou homomorphismes de
  Modules, ainsi qu'aux propriétés du type considéré dans (ii).

  \item Comme nous le verrons dans la suite de ce paragraphe et dans le courant
  du reste du chap.~IV, la plupart des propriétés $\mathbf P$ vérifiant la
  condition (9.2.1, (i)) vérifient aussi (9.2.1, (ii)). Comme exceptions
  possibles, notons la propriété d'être \emph{projectif}, ou
  \emph{quasi-projectif}, ou \emph{affine}, ou \emph{quasi-affine} sur le
  corps de base (pour un préschéma algébrique)~; nous verrons (9.6.2) que ces
  propriétés sont \emph{ind-constructibles}, mais nous donnerons plus tard un
  exemple où $S$ est une partie ouverte non vide de $\operatorname{Spec}(\mathbf Z)$
  (ou une partie ouverte d'une courbe elliptique sur un corps fini) et où
  toutes les fibres $X_s$

  